\documentclass[10pt,aspectratio=169]{beamer}
\input{../../_en_preambulo_beamer}
% Comandos y macros estadísticas compartidas para las presentaciones Beamer en inglés (Metropolis)

% Conjuntos numéricos básicos
\newcommand{\R}{\mathbb{R}}
\newcommand{\N}{\mathbb{N}}
\newcommand{\Q}{\mathbb{Q}}
\newcommand{\Z}{\mathbb{Z}}
\newcommand{\set}[1]{\left\{ #1 \right\}}
\newcommand{\abs}[1]{\left|#1\right|}
\newcommand{\norm}[1]{\left\| #1 \right\|}

% Comandos estadísticos y probabilísticos del curso
\newcommand{\Var}{\operatorname{Var}}
\newcommand{\cov}{\operatorname{Cov}}
\newcommand{\corr}{\rho}
\newcommand{\comb}[2]{\begin{pmatrix} #1 \\ #2 \end{pmatrix}}
\newcommand{\s}{\sigma}
\newcommand{\particion}{\mathfrak{P}}
\newcommand{\card}[1]{\#\left( #1 \right)}
\newcommand{\seq}[3]{\set{#1}_{#2}^{#3}}
\newcommand{\E}{\mathbb{E}}
\newcommand{\Prob}{\mathbb{P}}

% Entornos pedagógicos y cajas en inglés académico natural (soportando tanto nombres en inglés como en español)
\newenvironment{discussion}{%
  \begin{alertblock}{Discussion Question}%
}{%
  \end{alertblock}%
}
\newenvironment{discusion}{%
  \begin{alertblock}{Discussion Question}%
}{%
  \end{alertblock}%
}

\newenvironment{reflection}{%
  \begin{alertblock}{Classroom Reflection}%
}{%
  \end{alertblock}%
}
\newenvironment{reflexion}{%
  \begin{alertblock}{Classroom Reflection}%
}{%
  \end{alertblock}%
}

\newenvironment{paradox}{%
  \begin{alertblock}{Exploratory Paradox}%
}{%
  \end{alertblock}%
}
\newenvironment{paradoja}{%
  \begin{alertblock}{Exploratory Paradox}%
}{%
  \end{alertblock}%
}

\newenvironment{motivation}{%
  \begin{exampleblock}{Motivation: Data Science in Practice}%
}{%
  \end{exampleblock}%
}
\newenvironment{motivacion}{%
  \begin{exampleblock}{Motivation: Data Science in Practice}%
}{%
  \end{exampleblock}%
}

\newenvironment{quickchallenge}{%
  \begin{block}{Quick Team Challenge}%
}{%
  \end{block}%
}
\newenvironment{retoclase}{%
  \begin{block}{Quick Team Challenge}%
}{%
  \end{block}%
}


\title[Confidence Intervals]{Section 6.2: Confidence Intervals}
\subtitle{Confidence levels, tails, and $p$-values}
\author[J. Castillo Colmenares]{Juliho Castillo Colmenares}
\institute{Tecnológico de Monterrey}
\date{}

\begin{document}
\begin{frame}[plain]\vspace{-0.8cm}\titlepage\end{frame}

\begin{frame}{Roadmap --- Chapter 6}
  \footnotesize
  \begin{itemize}\itemsep=0.08cm
    \item \textbf{06.01} Point estimation
    \item \textbf{06.02} Confidence intervals \emph{(today)}
    \item \textbf{06.04--06.07} Means, proportions, variances, and sample size
  \end{itemize}
  \begin{block}{Learning objective}
    Build intervals for a mean, distinguish $Z$ from $t$, and connect confidence,
    significance, and $p$-values.
  \end{block}
\end{frame}

\begin{frame}{Motivation: a plausible range}
  \footnotesize
  A point estimator does not display all uncertainty. A confidence interval gives
  a range of values compatible with the data.
  \begin{alertblock}{Guiding question}
    How wide must the range be for the chosen confidence level?
  \end{alertblock}
\end{frame}

\begin{frame}{Confidence and $p$-value}
  \footnotesize
  For confidence level $p$, the reading defines significance as
  \[
    \beta=1-p.
  \]
  A statistic's $p$-value is written as
  \[
    p\text{-value}=P(X>Z)\quad\text{or}\quad P(X>t).
  \]
\end{frame}

\begin{frame}{Criterion and tails}
  \footnotesize
  \begin{itemize}\itemsep=0.12cm
    \item If $p\text{-value}>\beta$, do not reject the null hypothesis.
    \item If $p\text{-value}<\beta$, favor the alternative.
  \end{itemize}
  Normal symmetry gives left-tailed, right-tailed, and two-tailed tests.
\end{frame}

\begin{frame}{Interval for $\mu$ with known $\sigma$}
  \footnotesize
  \[
    \bar X-Z_{\alpha/2}\frac{\sigma}{\sqrt n}
    \le\mu\le
    \bar X+Z_{\alpha/2}\frac{\sigma}{\sqrt n}.
  \]
  The critical value comes from the standard normal distribution.
\end{frame}

\begin{frame}{Interval for $\mu$ with unknown $\sigma$}
  \footnotesize
  \[
    \bar X-t_{\alpha/2,n-1}\frac{S}{\sqrt n}
    \le\mu\le
    \bar X+t_{\alpha/2,n-1}\frac{S}{\sqrt n}.
  \]
  Replacing $\sigma$ by $S$ changes the reference distribution to $t$.
\end{frame}

\begin{frame}{Common structure}
  \footnotesize
  \[
    \text{point estimator}\ \pm\ (\text{critical value})
    (\text{standard error}).
  \]
  The same form recurs for means, proportions, and variances.
\end{frame}

\begin{frame}{Computational bridge 1/4: margin of error}
  \footnotesize
  \[
    E=(\text{critical value})\times SE.
  \]
  Increasing confidence increases the critical value and widens the interval.
\end{frame}

\begin{frame}{Computational bridge 2/4: Z versus t}
  \footnotesize
  \begin{tabular}{lll}
    \hline
    Case & Standard error & Reference \\
    \hline
    $\sigma$ known & $\sigma/\sqrt n$ & $Z$ \\
    $\sigma$ unknown & $S/\sqrt n$ & $t_{n-1}$ \\
    \hline
  \end{tabular}
\end{frame}

\begin{frame}{Computational bridge 3/4: tails}
  \footnotesize
  A two-sided test splits $\alpha$ between both tails. A one-sided test places
  the full critical region at one end.
  \begin{block}{Consequence}
    The tail type must be fixed before interpreting the $p$-value.
  \end{block}
\end{frame}

\begin{frame}{Computational bridge 4/4: interpretation}
  \footnotesize
  A 95\% CI uses $\alpha=0.05$ and $Z_{0.025}$ or $t_{0.025,n-1}$.
  \begin{alertblock}{Careful wording}
    The interval describes a repeated procedure; it does not make a fixed
    parameter random after the data are observed.
  \end{alertblock}
\end{frame}

\begin{frame}{Solved example 1/4 --- bottles and $Z$}
  \footnotesize
  A sample of $n=16$ bottles has $\bar X=500.4$ ml and known $\sigma=2$ ml.
  Construct the 95\% CI for $\mu$.
\end{frame}

\begin{frame}{Solved example 1/4 --- solution}
  \footnotesize
  \[
    SE=2/\sqrt{16}=0.5,\qquad
    CI=500.4\pm1.96(0.5)=[499.42,501.38]\text{ ml}.
  \]
  The normal reference is used because $\sigma$ is known.
\end{frame}

\begin{frame}{Solved example 2/4 --- bottles and $t$}
  \footnotesize
  With the same data, suppose $\sigma$ is unknown and $S=2.3$ ml. Use
  $t_{15,0.025}\approx2.131$ for the 95\% CI.
\end{frame}

\begin{frame}{Solved example 2/4 --- solution}
  \footnotesize
  \[
    SE=2.3/\sqrt{16}=0.575,\qquad
    CI=500.4\pm2.131(0.575)=[499.18,501.63]\text{ ml}.
  \]
  The $t$ interval is wider because $\sigma$ was estimated.
\end{frame}

\begin{frame}{Reading application 3/4 --- right tail}
  \footnotesize
  For a claim of increase, the critical region lies in the right tail.
  \begin{block}{Statement}
    Compare the $p$-value with $\beta$ using the right-tail probability of the
    observed statistic.
  \end{block}
\end{frame}

\begin{frame}{Reading application 3/4 --- decision}
  \footnotesize
  If the $p$-value is below $\beta$, the evidence is in the critical region and
  $H_0$ is rejected; otherwise it is not rejected.
\end{frame}

\begin{frame}{Reading application 4/4 --- two tails}
  \footnotesize
  For a difference in either direction, a two-sided test allocates significance
  to both tails.
  \begin{block}{Statement}
    Read the interval as the set of values not rejected by the matching
    two-sided test.
  \end{block}
\end{frame}

\begin{frame}{Reading application 4/4 --- connection}
  \footnotesize
  A value $\mu_0$ outside the two-sided CI gives a $p$-value below $\alpha$;
  a value inside is not rejected by the procedure.
  \begin{exampleblock}{Result}
    A CI and a two-sided test are two readings of the same margin of error.
  \end{exampleblock}
\end{frame}

\begin{frame}{Synthesis of the section}
  \footnotesize
  \begin{itemize}\itemsep=0.12cm
    \item Confidence fixes $\alpha$ and the critical value.
    \item $Z$ uses known $\sigma$; $t$ uses $S$.
    \item Tail type sets the critical region.
    \item CI and $p$-value connect estimation and decisions.
  \end{itemize}
\end{frame}

\begin{frame}{Curricular transition}
  \footnotesize
  \begin{block}{Established}
    An interval combines an estimator, reference distribution, and standard
    error.
  \end{block}
  \begin{alertblock}{Next section}
    \textbf{06.04 Standard Errors}: compare precision formulas for means,
    differences, and proportions.
  \end{alertblock}
\end{frame}
\end{document}
